% ===================================================================
%  TAULA N.º 11 — La Ciutat i els seus monuments. — L'Incendi.
%  Traduction 2026 d'apres texto/io, controlee sur texto/fr et
%  texto/es. Consigne : voir texto/ca/10-taula-01.tex.
% ===================================================================

%%K t11-apar-1 apar
\VUsaut{2.0mm}
\VUcentre{13.2pt}{90}{TERCERA SÈRIE}
\VUsaut{3.0mm}
\VUfilet{16mm}
\VUsaut{5.6mm}
\VUtitre{15.0pt}{60}{TAULA N.º 11}
\VUsaut{1.4mm}
\VUpk{11.4pt}{40}{La Ciutat i els seus monuments. --- L'Incendi.}
\VUsaut{2.2mm}

%%K t11-tit-1 sub
\VUpk{10.2pt}{40}{Els afores.}
\VUsaut{3.0mm}

%%K t11-c1-01-1 p
1. --- Visc en una gran ciutat situada a cent quilòmetres de l'oceà i que,
tanmateix, és un \VUgras{port} \textsuperscript{(1)} de primer ordre. El riu que la
banya té quatre-cents metres d'ample i forma una rada molt bella.

%%K t11-c1-02-1 p
2. --- Tota la part de la ciutat que és vora els \VUgras{molls}
\textsuperscript{(2)} té un aspecte completament marítim i industrial i forma un
\VUgras{raval} \textsuperscript{(3)} habitat només per mariners i obrers. Aquests
treballen a les \VUgras{fàbriques} \textsuperscript{(4)} i a la \VUgras{fàbrica de
gas} \textsuperscript{(5)}, la \VUgras{xemeneia} \textsuperscript{(6)} alta i el
\VUgras{gasòmetre} de la qual es veuen de molt lluny.

%%K t11-c1-03-1 p
3. --- A l'Edat Mitjana la ciutat estava fortificada, i encara es troben algunes
restes de les seves muralles, en particular la \VUgras{porta}
\textsuperscript{(8)} del vell \VUgras{castell fort} \textsuperscript{(7)},
flanquejat per les seves dues \VUgras{torres} \textsuperscript{(9)}, que avui
serveixen de \VUgras{presó}.

%%K t11-c1-04-1 p
4. --- En tot aquest \VUgras{barri}, d'aspecte molt antic, un sol edifici és
relativament modern: la \VUgras{duana} \textsuperscript{(10)}. És entre el moll i
el petit \VUgras{carrer} \textsuperscript{(11)} que porta a la vella \VUgras{casa
de la vila} \textsuperscript{(12)}. Aquest darrer monument es reconeix fàcilment
pel gegantí \VUgras{campanar} \textsuperscript{(13)} que domina tota la ciutat
vella.

%%K t11-c1-05-1 p
5. --- A l'altre costat del carrer \VUgras{s'alça} un edifici construït en el
mateix estil que la duana: és la \VUgras{borsa} \textsuperscript{(14)}. Una de les
seves façanes dóna a una placeta on hi ha el \VUgras{govern civil}
\textsuperscript{(15)}. La nostra ciutat, en efecte, sense ser una capital, és un
important cap de província.

%%K t11-tit-2 sub
\VUsaut{5.0mm}
\VUpk{10.2pt}{40}{La Part alta de la Ciutat.}
\VUsaut{3.4mm}

%%K t11-c1-06-1 p
6. --- La guarnició, força nombrosa, s'allotja en vastes \VUgras{casernes}
\textsuperscript{(16)} molt salubres; s'estenen fins a la reixa del
\VUgras{parc municipal} \textsuperscript{(17)}. El \VUgras{bulevard}
\textsuperscript{(19)} que porta a aquell parc passa per darrere de la \VUgras{caixa
d'estalvis} \textsuperscript{(18)} i desemboca a la plaça de la \VUgras{catedral}
\textsuperscript{(20)}.

%%K t11-c1-07-1 p
7. --- Tots aquests monuments s'esglaonen en amfiteatre damunt un turó on és també
el nostre ampli \VUgras{institut} \textsuperscript{(21)}.

%%K t11-c1-07-2 p
Si es baixa per una via una mica costeruda que s'ajunta amb el Carrer Major,
s'arriba al \VUgras{palau de justícia} \textsuperscript{(22)}, davant del qual
s'estén un agradable \VUgras{jardí públic} \textsuperscript{(26)}. Aquesta
\VUgras{plaça enjardinada} no és gaire gran, però posa al mig de la ciutat un
oasi de verdor i de frescor. Per això la freqüenten molt els habitants, als quals
els agrada de venir a seure sota el tendal del \VUgras{cafè}
\textsuperscript{(25)} per escoltar la banda militar que toca els dijous i els
diumenges al \VUgras{quiosc} \textsuperscript{(27)} posat entre les gespes i els
massissos.

%%K t11-c1-08-1 p
8. --- En una d'aquelles gespes s'alça una \VUgras{estàtua}
\textsuperscript{(28)} de bronze, monument erigit a la memòria d'un dels nostres
conciutadans que va prestar grans serveis a la seva ciutat natal.

%%K t11-tit-3 sub
\VUsaut{4.0mm}
\VUpk{10.2pt}{40}{Els transports.}
\VUsaut{3.4mm}

%%K t11-c1-09-1 p
9. --- La nostra ciutat encara no està enllumenada amb llum elèctrica; només té
\VUgras{fanals de gas} \textsuperscript{(29)}. Però la \VUgras{premsa local} fa
campanya perquè aquell sistema d'enllumenat es millori, \VUgras{igual que} ja
s'ha perfeccionat el \VUgras{sistema de tracció} dels tramvies. \VUgras{Els}
vells cotxes \VUgras{tirats} per \VUgras{cavalls} han estat substituïts per
pràctics \VUgras{tramvies elèctrics} \textsuperscript{(31)}, que porten ràpidament
els viatgers d'un extrem de \VUgras{la ciutat a l'altre} per un preu molt
\VUgras{barat}.

%%K t11-c1-10-1 p
10. --- Vora el palau de justícia hi ha el \VUgras{banc} \textsuperscript{(32)}, on
es descompten els valors comercials. Els \VUgras{ordenances del banc}
\textsuperscript{(33)} van a domicili a fer els cobraments.

%%K t11-tit-4 sub
\VUsaut{4.0mm}
\VUpk{10.2pt}{40}{Art i Instrucció.}
\VUsaut{3.4mm}

%%K t11-c1-11-1 p
11. --- El bell edifici que hi ha al costat és el \VUgras{museu}. Conté alhora la
\VUgras{biblioteca} \textsuperscript{(34)} de la ciutat, belles \VUgras{col·leccions
de ciències naturals} \textsuperscript{(35)} (museu d'història natural) i una
graciosa \VUgras{galeria de pintura} \textsuperscript{(36)} que constitueix una
esplèndida \VUgras{exposició} permanent.

%%K t11-c1-11-2 p
Està obert al públic dues vegades per setmana, però els forasters poden
visitar-lo qualsevol dia donant una propina al \VUgras{guarda}
\textsuperscript{(37)}. Els \VUgras{pintors} \textsuperscript{(38)} hi vénen a
treballar. Se'ls veu \VUgras{davant} del seu cavallet, amb la \VUgras{paleta} a
la mà, copiant els quadres cèlebres.

%%K t11-c1-12-1 p
12. --- A l'altura, darrere del museu, es distingeix la \VUgras{cúpula}
\textsuperscript{(40)} de l'\VUgras{hospital} \textsuperscript{(39)} i els edificis
de l'\VUgras{escola normal} \textsuperscript{(41)} en què es formaven els mestres
de les nostres escoles municipals. A la plaça del Teatre hi ha una \VUgras{oficina
de correus} \textsuperscript{(43)} instal·lada en una \VUgras{casa particular}
\textsuperscript{(42)}.

%%K t11-tit-5 sub
\VUsaut{5.0mm}
\VUpk{11.4pt}{40}{L'Incendi.}
\VUsaut{3.0mm}

%%K t11-tit-6 sub
\VUpk{10.2pt}{40}{El crit de «Foc!».}
\VUsaut{3.4mm}

%%K t11-c2-01-1 p
1. --- Una notícia alarmant s'escampa per la ciutat: acaba de declarar-se un
incendi al teatre! Quan arribo a la \VUgras{plaça} \textsuperscript{(96)}, la
multitud ja s'atapeeix a la \VUgras{vorera} \textsuperscript{(46)} i a la
\VUgras{calçada} \textsuperscript{(47)}. Els \VUgras{empleats de correus}
\textsuperscript{(44)} i els \VUgras{carters} \textsuperscript{(45)} surten de
l'oficina. Un \VUgras{conductor de tramvia} \textsuperscript{(48)} ha hagut
d'aturar el seu tramvia per deixar passar una \VUgras{ambulància}
\textsuperscript{(50)} municipal que arriba amb un \VUgras{cirurgià}
\textsuperscript{(52)} i un \VUgras{infermer} \textsuperscript{(51)} per atendre un
\VUgras{ferit} \textsuperscript{(53)} que han portat en una \VUgras{llitera}
\textsuperscript{(54)} fins a prop del \VUgras{quiosc de diaris}
\textsuperscript{(55)}.

%%K t11-c2-02-1 p
2. --- Una companyia d'infanteria que tornava de l'exercici s'ha aturat per
mantenir l'ordre amb els \VUgras{guàrdies municipals} \textsuperscript{(61)} a
cavall. Els \VUgras{genets}, els cavalls dels quals s'encabriten, ho fan millor
que els \VUgras{soldats} \textsuperscript{(57)} o els \VUgras{gendarmes}
\textsuperscript{(63)} per fer recular els \VUgras{curiosos}
\textsuperscript{(56)}. Altrament, aquests darrers estan molt ocupats. Un d'ells
indica als \VUgras{agents de policia} \textsuperscript{(64)} on cal portar una
altra \VUgras{víctima} \textsuperscript{(65)}, un valent bomber que ha caigut d'una
escala.

%%K t11-c2-03-1 p
3. --- L'\VUgras{oficial} \textsuperscript{(66)} precisa el lloc on ha de
col·locar-se la \VUgras{bomba de vapor} \textsuperscript{(67)} que va arribant.
Mentre espera aquella potent màquina, ha fet muntar una \VUgras{bomba de mà}
\textsuperscript{(68)}, la llarga \VUgras{mànega} \textsuperscript{(69)} de la qual
llança torrents d'aigua damunt el foc. Sis bombers la manegen mentre el seu
company, que sosté la \VUgras{llança} \textsuperscript{(70)}, dirigeix el
\VUgras{raig} \textsuperscript{(71)} al centre de l'incendi.

%%K t11-c2-04-1 p
4. --- A l'altre costat del teatre s'ha recolzat la gran \VUgras{escala}
\textsuperscript{(72)}. Permet als bombers d'acostar-se a les flames que broten
entre remolins de \VUgras{fum} \textsuperscript{(75)} negre.

%%K t11-tit-7 sub
\VUsaut{5.0mm}
\VUpk{10.2pt}{40}{Actes de valor.}
\VUsaut{3.4mm}

%%K t11-c2-05-1 p
5. --- L'\VUgras{incendi} \textsuperscript{(74)} va néixer al vestíbul del
\VUgras{teatre} \textsuperscript{(73)}, mentre s'assajava l'\VUgras{òpera} la
primera \VUgras{representació} de la qual havia de tenir lloc demà. Per sort
només hi havia uns pocs \VUgras{espectadors} \textsuperscript{(76)} a la sala. Els
pobres s'han refugiat al \VUgras{balcó} \textsuperscript{(77)}, on valents
\VUgras{bombers} \textsuperscript{(86)} acudeixen a socórrer-los amb una escala i
cordes.

%%K t11-c2-06-1 p
6. --- Sota el \VUgras{pòrtic} \textsuperscript{(78)}, on el \VUgras{cartell}
\textsuperscript{(79)} anuncia la representació, es veu fugir els \VUgras{músics}
\textsuperscript{(81)} de l'\VUgras{orquestra} \textsuperscript{(80)} enduent-se
els seus instruments: \VUgras{violí} \textsuperscript{(81)}, \VUgras{flauta}
\textsuperscript{(82)}, \VUgras{violoncel} \textsuperscript{(83)},
\VUgras{trombó} \textsuperscript{(84)}, \VUgras{tambor}
\textsuperscript{(85)}, etc. S'intenta salvar part del \VUgras{mobiliari}
\textsuperscript{(87)}.

%%K t11-c2-07-1 p
7. --- Els \VUgras{actors} \textsuperscript{(89)} i les \VUgras{actrius}
\textsuperscript{(88)}, \VUgras{cantants} ells i elles, han hagut de fugir amb el
seu \VUgras{vestit} \textsuperscript{(90)} de teatre. El tenor explica a un
\VUgras{periodista} \textsuperscript{(92)} com ha ocorregut. \VUgras{El reporter}
va acudir corrents des del primer moment amb les autoritats: el
\VUgras{governador} \textsuperscript{(91)}, l'\VUgras{alcalde}
\textsuperscript{(93)}, el \VUgras{comissari de policia} \textsuperscript{(94)} i
un \VUgras{jutge} \textsuperscript{(95)}, que obriran una investigació per
determinar les responsabilitats.
\VUsaut{6.0mm}
\VUfilet{22mm}
